% !TeX root = ../../Book5.tex
% 纲要标题：### 13.2 Cartan 判据
% 纲要位置：工作记录/计划纲要.md，第 4149 行。

\section{Cartan 判据}
\label{b5:sec:1302}

% 纲要范围：
%
% * 可解性的迹判据
% * 半单性的非退化判据
% * 与 Engel 定理的关系
%

迹型给出的二次条件足以判别可解性。不过，复数特征值的平方可能相消，不能从 \(\operatorname{tr}(x^2)=0\) 直接推出 \(x=0\)。证明中将另构造一个算子，把所用的迹写成特征值绝对值平方之和。

\subsection{可解性的迹判据}
\label{b5:subsec:1302:01}

\begin{lemma}[谱插值]\label{b5:lem:cartan-spectral-interpolation}
设 \(V\) 为有限维复向量空间，\(x\in\operatorname{End}_{\mathbb C}(V)\)，\(V_\lambda\) 表示 \(x\) 对应于特征值 \(\lambda\) 的广义特征空间。按这些空间写 \(x=S+N\)，其中 \(S\) 在 \(V_\lambda\) 上为 \(\lambda I\)，\(N\) 幂零。令 \(S^\#\) 在 \(V_\lambda\) 上为 \(\overline\lambda I\)。则存在复多项式 \(p\)，满足 \(p(0)=0\)，并且
\[
\operatorname{ad}S^\#=p(\operatorname{ad}x)
\quad\text{在 }\operatorname{End}_{\mathbb C}(V)\text{ 上}.
\]
\end{lemma}
\begin{proof}
在 \(\operatorname{Hom}(V_\mu,V_\lambda)\) 上，
\(\operatorname{ad}x=(\lambda-\mu)I+\operatorname{ad}N\)，其中后项幂零；而 \(\operatorname{ad}S^\#\) 是标量 \(\overline{\lambda-\mu}\)。取一个对所有这些子空间都足够大的整数 \(r\)。对有限多个不同的差值 \(\delta=\lambda-\mu\)，多项式中国剩余定理给出
\[
p(t)\equiv\overline\delta\pmod{(t-\delta)^r}.
\]
当 \(\delta=0\) 时，这要求 \(p(t)\equiv0\pmod{t^r}\)，特别有 \(p(0)=0\)。代入每个子空间上的算子即可；这些 Hom 空间的直和就是全部自同态空间。
\end{proof}

\term[Cartan-kejie-panju]{Cartan 可解性判据}[Cartan's solvability criterion]如下。

\begin{theorem}[Cartan 可解性判据]\label{b5:thm:cartan-solvability}
设 \(L\subseteq\mathfrak{gl}(V)\) 是有限维复向量空间 \(V\) 上的李子代数。则 \(L\) 可解，当且仅当
\[
\operatorname{tr}(xy)=0
\quad\Bigl(x\in[L,L],\ y\in L\Bigr).
\]
因而有限维复李代数 \(\mathfrak g\) 可解，当且仅当
\(\kappa_{\mathfrak g}\Bigl([\mathfrak g,\mathfrak g],\mathfrak g\Bigr)=0\)。
\end{theorem}
\begin{proof}
若 \(L\) 可解，\cref{b5:prop:solvable-derived-nilpotent}已经证明迹条件。

反之，固定 \(x\in[L,L]\)，采用\cref{b5:lem:cartan-spectral-interpolation}的 \(S^\#\)。因为 \(\operatorname{ad}x\) 把 \(L\) 送入理想 \([L,L]\)，而 \(p\) 没有常数项，所以
\[
[S^\#,L]=p(\operatorname{ad}x)(L)\subseteq[L,L].
\]
把 \(x\) 写成有限和 \(\sum_j[a_j,b_j]\)，其中 \(a_j,b_j\in L\)。迹的循环性和假设给出
\[
\operatorname{tr}(S^\#x)
=\sum_j\operatorname{tr}\Bigl([S^\#,a_j]b_j\Bigr)=0.
\]
另一方面，在 \(x\) 的 Jordan 基下，
\[
\operatorname{tr}(S^\#x)
=\sum_\lambda(\dim V_\lambda)|\lambda|^2.
\]
每一项均非负，故每个特征值都为零，\(x\) 幂零。于是 \([L,L]\) 中的每个矩阵都是幂零算子。\cref{b5:thm:engel-common-zero}使其同时严格上三角化，故它是幂零李代数，特别可解；所以 \(L\) 可解。

最后，对抽象李代数取 \(L=\operatorname{ad}\mathfrak g\)，矩阵迹条件正好是陈述中的 Killing 型条件。若该条件成立，已经证明 \(L\) 可解，而
\(\mathfrak g/Z(\mathfrak g)\cong L\)。中心交换，由\cref{b5:prop:lie-solvable-extension}，\(\mathfrak g\) 可解。逆向对伴随表示使用已证的必要性。
\end{proof}

\subsection{半单性的非退化判据}
\label{b5:subsec:1302:02}

\term[Cartan-bandan-panju]{Cartan 半单性判据}[Cartan's semisimplicity criterion]给出一个更简洁的检验。

\begin{theorem}[Killing 型的非退化判据]\label{b5:thm:killing-semisimple}
有限维复李代数 \(\mathfrak g\) 半单，当且仅当其 Killing 型非退化。
\end{theorem}
\begin{proof}
设 \(\mathfrak g\) 半单，令 \(R=\mathfrak g^\perp\) 为 Killing 型的根。\cref{b5:prop:invariant-trace-form}说明 \(R\) 是理想。\cref{b5:lem:killing-restriction-ideal}给出 \(\kappa_R=\kappa_{\mathfrak g}|_{R\times R}=0\)，故\cref{b5:thm:cartan-solvability}说明 \(R\) 可解。半单性迫使 \(R=0\)。

反之，设 Killing 型非退化。若有非零可解理想 \(I\)，其导出列的最后一个非零项 \(A\) 是非零交换理想：各导出项均为 \(\mathfrak g\) 的理想，且 \([A,A]=0\)。对 \(a\in A\)、\(x\in\mathfrak g\)，算子 \(T=\operatorname{ad}a\operatorname{ad}x\) 把 \(\mathfrak g\) 送入 \(A\)，又把 \(A\) 送到零，所以 \(T^2=0\)。于是 \(\kappa_{\mathfrak g}(a,x)=\operatorname{tr}T=0\)。这使 \(A\) 包含于型的根，与非退化性矛盾。因此不存在非零可解理想。
\end{proof}

\begin{corollary}\label{b5:cor:simple-killing-forms}
非零有限维复简单李代数 \(\mathfrak g\) 的对称不变双线性型均为 Killing 型的标量倍数，且 Killing 型非退化。
\end{corollary}
\begin{proof}
简单李代数半单，所以\cref{b5:thm:killing-semisimple}保证 Killing 型非退化；再用\cref{b5:prop:simple-invariant-form-unique}。
\end{proof}

\begin{proposition}\label{b5:prop:faithful-trace-nondegenerate}
设 \(\mathfrak g\) 是有限维复半单李代数，\(\rho:\mathfrak g\rightarrow\mathfrak{gl}(V)\) 是忠实有限维表示。则 \(\beta_\rho(x,y)=\operatorname{tr}\Bigl(\rho(x)\rho(y)\Bigr)\) 非退化。
\end{proposition}
\begin{proof}
型的根 \(R\) 是理想，且 \(\operatorname{tr}\Bigl(\rho(x)\rho(y)\Bigr)=0\) 对全部 \(x,y\in R\) 成立。对矩阵代数 \(\rho(R)\) 用\cref{b5:thm:cartan-solvability}，得 \(\rho(R)\) 可解。忠实性给出 \(R\cong\rho(R)\)，故 \(R\) 是可解理想，只能为零。
\end{proof}

\subsection{与 Engel 定理的关系}
\label{b5:subsec:1302:03}

证明的顺序值得区分。Engel 定理先由共同零向量识别幂零性；Lie 定理随后给出可解表示的上三角形；Cartan 判据再通过谱插值把迹条件还原为 Engel 定理的假设。半单性判据用的是这条已经建立的证明链，并未调用半单李代数表示的完全可约性。

\begin{example}\label{b5:exm:killing-affine-two}
设 \(\mathfrak g=\mathbb Ch+\mathbb Ce\)，\([h,e]=e\)。相对于 \(h,e\)，
\[
\operatorname{ad}h=\begin{pmatrix}0&0\\0&1\end{pmatrix},
\qquad
\operatorname{ad}e=\begin{pmatrix}0&0\\-1&0\end{pmatrix},
\qquad
[\kappa]_{h,e}=\begin{pmatrix}1&0\\0&0\end{pmatrix}.
\]
它可解而 Killing 型不恒为零；正确条件是
\(\kappa\Bigl([\mathfrak g,\mathfrak g],\mathfrak g\Bigr)=0\)，这里
\([\mathfrak g,\mathfrak g]=\mathbb Ce\) 恰好等于型的根。
\end{example}

特征零假设同样不可漏写。以特征 \(3\) 的域为例，
\(\mathfrak{sl}_2\) 仍简单：\(\operatorname{ad}h\) 在 \(e,h,f\) 上的三个特征值互不相同，任意非零理想可投影出一个基向量，再用三条标准括号得到全部基。但若特征为 \(2\)，这些步骤与 Killing 型公式都会退化，\(\mathfrak{sl}_2\) 不再半单。因此本章判据的适用范围始终是有限维复李代数，不能只看相同的矩阵记号就推广。
